\subsection{Equal liabilities can require different cash paths}
\label{sec:apple-settlement-paths}

The conditional repair leaves different obligations beneath the same annual
liability total. In the $q=2{,}000$ illustration, History 1 owes 1,082 for
capital purchases and 2,000 for the separate expense; History 2 owes 2,082
and 1,000 respectively. Each owes 3,082. Consider what paying those balances
would do, conditional on the economic and reporting assumptions needed for
the preceding construction. Assume payment at carrying value, sufficient
unrestricted cash, and no additional interest, foreign exchange,
remeasurement or tax effect. Capital settlement is classified as investing
cash and expense settlement as operating cash in this example. Neither
the obligation nor its settlement terms have been identified in Apple's
disclosures.

If both obligations are paid at the first subsequent date, each payment
debits its liability and credits cash. No new expense arises: the expense
or acquisition was recognised in the prior period. Both histories reduce
cash and liabilities by 3,082, leave assets other than cash unchanged, and
leave profit, OCI and equity unchanged. The indirect operating-cash
calculation includes the expense liability's decrease once; it excludes
capital settlement from the operating-liability adjustment. The
cash-flow statement nevertheless distinguishes the histories. Writing
each ordered pair as operating cash followed by investing cash, the
joint alternatives are
\[
(\mathrm{CFO},\mathrm{CFI})_{H_1}=(-2{,}000,-1{,}082),
\qquad
(\mathrm{CFO},\mathrm{CFI})_{H_2}=(-1{,}000,-2{,}082).
\]
Both pairs sum to $-3{,}082$. Selecting the smaller outflow from each
category would instead produce $(-1{,}000,-1{,}082)$, whose total is only
$-2{,}082$. No branch produces that pair. Selecting both larger outflows
would invent an extra 1,000 payment. The category intervals can describe
marginal possibilities, but their Cartesian product contains paths
excluded by the shared history.

Timing makes the distinction relevant even to total cash. Suppose capital
obligations are paid at the first date and expense obligations at the
second. History 1 then pays 1,082 followed by 2,000; History 2 pays 2,082
followed by 1,000. Relative to History 1, History 2 has 1,000 less cash
after the first date and the same cumulative cash after the second.
Paying the expense obligations first reverses the first-date difference.
An aggregate liability of 3,082 suffices for the narrow query of total
undiscounted cash after both payments, under these fixed assumptions.
It does not suffice for cash by date or by statement category. Retaining
liability kind, amount, settlement terms and joint branch identity is
therefore necessary for those queries.

The actual following quarter constrains a proposed continuation without
identifying it. Apple reports PP\&E payments of 2,940, operating cash of
29,935, investing cash of 9,792 and total cash growth of 356 in the quarter
ended 28 December 2024 \citep[printed p.~5]{apple2025q1}. If both
hypothetical obligations settled in that quarter, the remaining PP\&E
payments would have to be 1,858 under History 1 or 858 under History 2.
The rest of the company would have to contribute operating cash of
31,935 or 30,935, and investing cash of 10,874 or 11,874, respectively.
After unchanged financing cash of $-39{,}371$, either required remainder
has cash growth of 3,438; subtracting the hypothetical 3,082 settlement
recovers the reported 356. These are necessary conditions on unconstructed
remainders. Positive remaining payments and a matching cash identity do
not establish the underlying transactions or prove that either
continuation fits the complete quarterly record.

This calculation extends the conditional history into explicit journals
and statement effects while keeping actual-history admission unresolved.
If cash is insufficient, the model must retain the resulting funding or
default question; it cannot invent a loan to preserve the path. A later
depreciation or production query requires different information again:
asset cohorts, remaining lives, residual values, capacity and tax bases.
The settlement example supplies none of those quantities. A change in the
query must therefore reopen the state-sufficiency check, even when the
annual statements and the final cash total coincide.
